\subsection{Processing threads and buffers}\label{sec:threads}

Inside \eeye{} a producer thread captures, an effects thread runs the operators, and
a consumer thread publishes (Figure~\ref{fig:eeye}). The threads pass frames through
three ring buffers of 32 slots, and frames come from a pool of eight that is allocated
once and never grows, so that the pool size bounds the lag between capture and output
(eight frames, 267~ms at 30~frames per second). Each frame has three planar buffers:
the untouched capture (\emph{raw}), the working image, and a spare that lets two
consecutive neighbourhood operators alternate without copying.

The configuration is a JSON file that a watcher thread reloads within milliseconds of
any change (inotify). A parsed configuration is published by an atomic pointer swap,
and the effects thread takes one snapshot per frame, so that no frame is processed
with two configurations. After every load \eeye{} writes whether it accepted the file;
\hellion{}, which writes the file for the browser page, waits for that verdict and
puts the previous file back when the new one is refused. Once a second the consumer
assembles a health record (capture state, frame count, recording state and the image
statistics of Section~\ref{sec:autovis}); the stats thread writes it, so that file
input/output on the vehicle's SD card cannot stall the thread that publishes the
pilot's picture.

\begin{figure*}[!t]
  \centering
  \begin{adjustbox}{max width=\linewidth}
  \begin{tikzpicture}[node distance=5mm and 6mm]
    \node[ext] (cam) {camera or\\frame bus};
    \node[proc, right=of cam] (prod) {producer\\capture, decode};
    \node[buf, right=of prod] (rin) {ring\\\emph{in}};
    \node[proc, right=of rin, text width=34mm] (eff) {effects\\image statistics\\visibility correction\\operator chain};
    \node[buf, right=of eff] (rout) {ring\\\emph{out}};
    \node[proc, right=of rout] (cons) {consumer\\publish};
    \node[buf, below=9mm of eff] (rfree) {ring \emph{free}\\pool of 8 frames};
    \node[buf, right=11mm of cons, yshift=7mm] (bus) {frame bus\\\emph{pilot}, \emph{raw}};
    \node[ext, right=11mm of cons, yshift=-7mm] (out) {webcam device,\\raw recording};
    \node[proc, above=7mm of eff] (cfg) {config watcher};
    \node[ext, left=8mm of cfg] (file) {effects.json};
    \node[proc, above=7mm of cons] (stats) {stats\\health file};
    \draw[flow] (cam) -- (prod);
    \draw[flow] (prod) -- (rin);
    \draw[flow] (rin) -- (eff);
    \draw[flow] (eff) -- (rout);
    \draw[flow] (rout) -- (cons);
    \draw[flow] (cons.east) -- ++(3mm,0) |- (bus.west);
    \draw[flow] (cons.east) -- ++(3mm,0) |- (out.west);
    \draw[flow] (cons.south) |- (rfree.east);
    \draw[flow] (rfree.west) -| (prod.south);
    \draw[flow] (file) -- (cfg);
    \draw[flow, dashed] (cfg) -- node[lab, right] {one snapshot\\per frame} (eff);
    \draw[flow, dashed] (cons) -- node[lab, left] {once a second} (stats);
  \end{tikzpicture}
  \end{adjustbox}
  \caption{Threads and buffers of \eeye{}. Frames circulate from the free ring through
    capture, processing and publication and back; the pool of eight frames bounds the
    output lag. Configuration changes reach the effects thread as whole snapshots.
    The stats thread writes the health file, so that file input/output never runs on
    the consumer.}\label{fig:eeye}
\end{figure*}
